\documentclass[english,hebrew]{article}
\usepackage[HE8,T1]{fontenc}
\usepackage[utf8]{inputenc}
\setcounter{secnumdepth}{-2}
\setlength{\parindent}{0bp}
\usepackage{calc}
\usepackage{amsmath}
\usepackage{amsthm}
\usepackage{amssymb}

\makeatletter
\theoremstyle{definition}
\newtheorem*{problem*}{\protect\problemname}
\theoremstyle{definition}
\newtheorem*{sol*}{\protect\solutionname}

\@ifundefined{date}{}{\date{}}
\usepackage{graphicx}
\usepackage{geometry}
\newcommand{\limi}{\underset{n\rightarrow\infty}{\lim}}
\newcommand{\limxz}{\underset{x\rightarrow x_{0}}{\lim}}
\newcommand{\limfxz}{\underset{x\rightarrow x_{0}}{\lim}f(x)}
\newcommand{\limfxm}{\underset{x\rightarrow x_{0}^{-}}{\lim}f(x)}
\newcommand{\limfxp}{\underset{x\rightarrow x_{0}^{+}}{\lim}f(x)}
\newcommand{\limxm}{\underset{x\rightarrow x_{0}^{-}}{\lim}}
\newcommand{\limxp}{\underset{x\rightarrow x_{0}^{+}}{\lim}}
\newcommand{\sqrm}{M_{m\times m}(\mathbb{F})}
\newcommand{\seqa}{(a_{n})_{n=1}^{\infty}}
\newcommand{\seqx}{(x_{n})_{n=1}^{\infty}}
\newcommand{\funcdr}{f:D\rightarrow\mathbb{R}}
\newcommand{\der}{\limxz\frac{f(x)-f(x_{0})}{x-x_{0}}}
\usepackage[all]{pdfprivacy}

\makeatother

\usepackage{babel}
\addto\captionsenglish{\renewcommand{\problemname}{Problem}}
\addto\captionsenglish{\renewcommand{\solutionname}{Solution}}
\addto\captionshebrew{\renewcommand{\problemname}{בעיה}}
\addto\captionshebrew{\renewcommand{\solutionname}{פתרון}}
\providecommand{\problemname}{בעיה}
\providecommand{\solutionname}{פתרון}

\begin{document}
\newgeometry{top=1cm, bottom=2cm, left=1cm, right=1cm}
\title{חשבון אינפיניטסימלי {\beginL 2\endL} - לקט שאלות - שבוע {\beginL 1\endL}}
\author{נעם לב}

\maketitle
\noindent\fbox{\begin{minipage}[t]{1\columnwidth - 2\fboxsep - 2\fboxrule}%
\begin{center}
האוניברסיטה העברית \textbar{} חשבון אינפיניטסימלי ){\beginL 2\endL}(
- \L{$80132$} \textbar{} סמסטר{\small{} }{\beginL 2027\endL}{\small{}
}א' )תשפ\textquotedblleft ז א'( \textbar{} מרצה: דוקטור איב גודין
\textbar{} מתרגל: מר יואב כהן
\par\end{center}%
\end{minipage}}
\begin{center}
לקט השאלות מכיל שאלות ממקורות שונים.\\
מקור חלק מן השאלות הוא חומר הקורס, והזכויות עליו שייכות לאוניברסיטה
העברית ולסגל הקורס.\\
הלקט אינו רשמי ולא עבר אישור של הסגל. בפרט, אם ישנן טעויות, הן שלי
ולא של הסגל.
\par\end{center}

\medskip{}

\tableofcontents{}

\newpage{}

\part{פונקציות קדומות )אינטגרל לא מסוים(}

הנימוק \L{$\text{linearity}$} מתייחס ללינאריות האינטגרל הלא מסוים.
\begin{problem*}
חשבו את האינטגרל הלא מסוים \L{%
\[
\int\frac{1}{3x+4}dx
\]
}%
 עבור \L{$x>0$}.
\end{problem*}
\begin{sol*}
על מנת לפתור את הבעיה, נשתמש בכתיב לייבניץ לנגזרת ובמשפט אינטגרציה
בהצבה.

נזהה שהאינטגרל הלא מסוים \L{$\int\frac{1}{x}dx$} הוא קל לחישוב ודומה
למבנה האינטגרל הלא מסוים שלנו, על כן נגדיר \L{%
\[
u=3x+4
\]
}%
בשביל להציב את \L{$u$} באינטגרל הלא מסוים שלנו, נרצה לחשב את האינטגרל
הלא מסוים לפי \L{$u$} ולא לפי \L{$x$}.

לשם כך, נצטרך לבטא את \L{$dx$} בעזרת \L{$du$}:\L{%
\[
\frac{du}{dx}=u'=3
\]
}%
ועל כן\L{%
\[
dx=\frac{1}{3}du
\]
}%
לכן נוכל לכתוב\L{%
\[
\int\frac{1}{3x+4}dx=\int\frac{1}{u}\ \frac{1}{3}du\underset{\text{linearity}}{=}\frac{1}{3}\int\frac{1}{u}\ du=\frac{1}{3}\ln u+C=\frac{1}{3}\ln\left(3x+4\right)+C
\]
}%
\end{sol*}
\begin{problem*}
חשבו את האינטגרל הלא מסוים \L{%
\[
\int\frac{1}{4x^{2}+12x+10}dx
\]
}%
\end{problem*}
\begin{sol*}
על מנת לפתור את הבעיה, נשתמש בכתיב לייבניץ לנגזרת ובמשפט אינטגרציה
בהצבה.

נזהה שהאינטגרל הלא מסוים \L{$\int\frac{1}{x^{2}+1}dx$} הוא קל לחישוב
ודומה למבנה האינטגרל הלא מסוים שלנו - על כן ננסה \textquotedblleft לייצר\textquotedblright{}
אותו.\\
\noindent\fbox{\begin{minipage}[t]{1\columnwidth - 2\fboxsep - 2\fboxrule}%
\begin{center}
טיוטה:\L{%
\[
mn=4\cdot9=36
\]
}%
\L{%
\[
m+n=12
\]
}%
\L{%
\[
m=n=6
\]
}%
\L{%
\[
4x^{2}+12x+9=\left(2x+3\right)^{2}
\]
}%
\par\end{center}%
\end{minipage}}

\L{%
\[
\frac{1}{4x^{2}+12x+10}=\frac{1}{4x^{2}+12x+9+1}=\frac{1}{\left(2x+3\right)^{2}+1}
\]
}%

ולכן נגדיר \L{%
\[
u=2x+3
\]
}%
\L{%
\[
\frac{du}{dx}=2\Rightarrow dx=\frac{1}{2}du
\]
}%
\L{%
\[
\therefore\int\frac{1}{4x^{2}+12x+10}dx\underset{\text{linearity}}{=}\frac{1}{2}\int\frac{1}{u^{2}+1}\ du=\frac{1}{2}\arctan\left(u\right)+C=\frac{1}{2}\arctan\left(2x+3\right)+C
\]
}%
\end{sol*}
\begin{problem*}
חשבו את האינטגרל הלא מסוים\L{%
\[
\int\sin x\cos x\ dx
\]
}%
\end{problem*}
\begin{sol*}
נשים לב שמתקיים\L{%
\[
\int\sin x\cos x\ dx=\int\sin x\sin'x\ dx
\]
}%
נציב \L{$u=\sin x$} ונקבל\L{%
\[
\frac{du}{dx}=\sin'x\Rightarrow du=\sin'x\ dx
\]
}%
\L{%
\[
\therefore\int\sin x\sin'x\ dx=\int u\ du=\frac{u^{2}}{2}+C=\frac{\sin^{2}x}{2}+C
\]
}%
\end{sol*}
\begin{problem*}
חשבו את האינטגרל הלא מסוים\L{%
\[
\int\sin^{6}(x)\cos^{5}(x)\ dx
\]
}%
\end{problem*}
\begin{sol*}
יעלה בהמשך
\end{sol*}
\begin{problem*}
חשבו את האינטגרל הלא מסוים\L{%
\[
\int2xe^{x^{2}}\ dx
\]
}%
\end{problem*}
\begin{sol*}
יעלה בהמשך
\end{sol*}
\begin{problem*}
חשבו את האינטגרל הלא מסוים\L{%
\[
\int e^{x}\sin x\ dx
\]
}%
\end{problem*}
\begin{sol*}
יעלה בהמשך
\end{sol*}
\begin{problem*}
חשבו את האינטגרל הלא מסוים\L{%
\[
\int\frac{1}{1+x^{2}}\ dx
\]
}%
\end{problem*}
\begin{sol*}
יעלה בהמשך
\end{sol*}
\begin{problem*}
חשבו את האינטגרל הלא מסוים\L{%
\[
\int\frac{x}{1+x^{2}}\ dx
\]
}%
\end{problem*}
\begin{sol*}
יעלה בהמשך
\end{sol*}

\end{document}
